% GENERATED FILE. Edit canonical lessons, then run npm run generate:books.
% canonical-source-hash: fnv1a64:08d1eaf1455502e8
% canonical-lessons: GU-C479-madu, GU-C479-lakvo, GU-C479-acki, GU-C479-gath, GU-C479-patharivas

\chapter{Taken Ill}
\label{ch:gu-taken-ill}
\hlchaptermodality{479}

\begin{chapteropening}
\textbf{By the end of this chapter:} \emph{I can say how serious an illness is.}

Complete the last lesson of chapter 479: I can say how serious an illness is.
\end{chapteropening}

\section[\texorpdfstring{\gu{માંદું} (mā̃dũ)}{માંદું (mā̃dũ)}]{\gu{માંદું} (mā̃dũ) — unwell; sick}

\label{lesson:GU-C479-madu}



\begin{quote}
Before the new one: say the Gujarati for to cough, then the Gujarati for constipation.
\end{quote}

\subsection*{You'll want to know: \gu{માંદું}}
\textbf{\gu{માંદું}} — \emph{mā̃dũ} — \textquotedblleft{}unwell; sick\textquotedblright{}.

\begin{quote}
\gu{મીરા} \gu{ત્રણ} \gu{દિવસથી} \gu{માંદી} \gu{છે}.
\end{quote}

\begin{quote}
\emph{Mīrā traṇ divasthī mā̃dī che.}
\end{quote}

\begin{quote}
Mira has been unwell for three days.
\end{quote}

\begin{cousinweb}
\textbf{\gu{માંદું}} (\emph{mā̃dũ}) goes back to a reconstructed \emph{*matthara-}, \textquotedblleft{}defective\textquotedblright{}, beside Sanskrit \emph{manthara}, \textquotedblleft{}slow, dull\textquotedblright{}; CDIAL lists the Gujarati word, \textquotedblleft{}sick\textquotedblright{} (CDIAL 9754).
\end{cousinweb}

\begin{grammarlens}[title={What you've built}]
One more word for describing what you lived and felt.
\end{grammarlens}

\subsection*{Your turn}
\begin{itemize}
\raggedright
  \item \emph{Say it:} \emph{mā̃dũ}
  \item \emph{Say it:} \emph{mā̃dũ}, once more
  \item \emph{Say it:} \emph{Mīrā traṇ divasthī mā̃dī che.} — Mira has been unwell for three days.
  \item \emph{Recall:} say the Gujarati for to cough, then the Gujarati for constipation, then say \emph{mā̃dũ} again
\end{itemize}

\begin{tcolorbox}[breakable,colback=teal!4,colframe=teal!35!black,title={Before you move on}]
What does \gu{માંદું} mean? (\textquotedblleft{}Unwell; sick\textquotedblright{}.) Say it once more.
\end{tcolorbox}
\section[\texorpdfstring{\gu{લકવો} (lakvo)}{લકવો (lakvo)}]{\gu{લકવો} (lakvo) — paralysis}

\label{lesson:GU-C479-lakvo}



\begin{quote}
Before the new one: say the Gujarati for constipation, then the Gujarati for unwell.
\end{quote}

\subsection*{You'll want to know: \gu{લકવો}}
\textbf{\gu{લકવો}} — \emph{lakvo} — \textquotedblleft{}paralysis\textquotedblright{}.

\begin{quote}
\gu{ગયા} \gu{વરસે} \gu{દાદાને} \gu{લકવો} \gu{થયો}.
\end{quote}

\begin{quote}
\emph{gayā varse dādāne lakvo thayo.}
\end{quote}

\begin{quote}
Last year my grandfather was paralysed by a stroke.
\end{quote}

A masculine noun.

\begin{cousinweb}
\textbf{\gu{લકવો}} (\emph{lakvo}) is usually traced to Arabic \emph{laqwa}, \textquotedblleft{}facial palsy\textquotedblright{}, through Persian.
\end{cousinweb}

\begin{grammarlens}[title={What you've built}]
One more word for describing what you lived and felt.
\end{grammarlens}

\subsection*{Your turn}
\begin{itemize}
\raggedright
  \item \emph{Say it:} \emph{lakvo}
  \item \emph{Say it:} \emph{lakvo}, once more
  \item \emph{Say it:} \emph{gayā varse dādāne lakvo thayo.} — Last year my grandfather was paralysed by a stroke.
  \item \emph{Recall:} say the Gujarati for constipation, then the Gujarati for unwell, then say \emph{lakvo} again
\end{itemize}

\begin{tcolorbox}[breakable,colback=teal!4,colframe=teal!35!black,title={Before you move on}]
What does \gu{લકવો} mean? (\textquotedblleft{}Paralysis\textquotedblright{}.) Say it once more.
\end{tcolorbox}
\section[\texorpdfstring{\gu{આંચકી} (ā̃ckī)}{આંચકી (ā̃ckī)}]{\gu{આંચકી} (ā̃ckī) — a fit; a convulsion}

\label{lesson:GU-C479-acki}



\begin{quote}
Before the new one: say the Gujarati for unwell, then the Gujarati for paralysis.
\end{quote}

\subsection*{You'll want to know: \gu{આંચકી}}
\textbf{\gu{આંચકી}} — \emph{ā̃ckī} — \textquotedblleft{}a fit; a convulsion\textquotedblright{}.

\begin{quote}
\gu{તાવમાં} \gu{બાળકને} \gu{આંચકી} \gu{આવી}.
\end{quote}

\begin{quote}
\emph{tāvmā̃ bāḷakne ā̃ckī āvī.}
\end{quote}

\begin{quote}
The child had a convulsion during the fever.
\end{quote}

A feminine noun.

\begin{cousinweb}
\textbf{\gu{આંચકી}} (\emph{ā̃ckī}) goes back to a reconstructed \emph{*acca-}, \textquotedblleft{}sudden shock\textquotedblright{}; CDIAL lists the Gujarati \emph{ācko}, \textquotedblleft{}shock, push\textquotedblright{} (CDIAL 141), and this word is built on it.
\end{cousinweb}

\begin{grammarlens}[title={What you've built}]
One more word for describing what you lived and felt.
\end{grammarlens}

\subsection*{Your turn}
\begin{itemize}
\raggedright
  \item \emph{Say it:} \emph{ā̃ckī}
  \item \emph{Say it:} \emph{ā̃ckī}, once more
  \item \emph{Say it:} \emph{tāvmā̃ bāḷakne ā̃ckī āvī.} — The child had a convulsion during the fever.
  \item \emph{Recall:} say the Gujarati for unwell, then the Gujarati for paralysis, then say \emph{ā̃ckī} again
\end{itemize}

\begin{tcolorbox}[breakable,colback=teal!4,colframe=teal!35!black,title={Before you move on}]
What does \gu{આંચકી} mean? (\textquotedblleft{}A fit; a convulsion\textquotedblright{}.) Say it once more.
\end{tcolorbox}
\section[\texorpdfstring{\gu{ગાંઠ} (gā̃ṭh)}{ગાંઠ (gā̃ṭh)}]{\gu{ગાંઠ} (gā̃ṭh) — a lump; a knot}

\label{lesson:GU-C479-gath}



\begin{quote}
Before the new one: say the Gujarati for paralysis, then the Gujarati for a convulsion.
\end{quote}

\subsection*{You'll want to know: \gu{ગાંઠ}}
\textbf{\gu{ગાંઠ}} — \emph{gā̃ṭh} — \textquotedblleft{}a lump; a knot\textquotedblright{}.

\begin{quote}
\gu{દાક્તરે} \gu{ગળામાં} \gu{નાની} \gu{ગાંઠ} \gu{જોઈ}.
\end{quote}

\begin{quote}
\emph{dāktare gaḷāmā̃ nānī gā̃ṭh joī.}
\end{quote}

\begin{quote}
The doctor found a small lump in the throat.
\end{quote}

A feminine noun.

\begin{cousinweb}
\textbf{\gu{ગાંઠ}} (\emph{gā̃ṭh}) is inherited from Sanskrit \emph{granthi}, \textquotedblleft{}knot, bunch, protuberance\textquotedblright{}; CDIAL lists the Gujarati word, \textquotedblleft{}knot\textquotedblright{} (CDIAL 4354).
\end{cousinweb}

\begin{grammarlens}[title={What you've built}]
One more word for describing what you lived and felt.
\end{grammarlens}

\subsection*{Your turn}
\begin{itemize}
\raggedright
  \item \emph{Say it:} \emph{gā̃ṭh}
  \item \emph{Say it:} \emph{gā̃ṭh}, once more
  \item \emph{Say it:} \emph{dāktare gaḷāmā̃ nānī gā̃ṭh joī.} — The doctor found a small lump in the throat.
  \item \emph{Recall:} say the Gujarati for paralysis, then the Gujarati for a convulsion, then say \emph{gā̃ṭh} again
\end{itemize}

\begin{tcolorbox}[breakable,colback=teal!4,colframe=teal!35!black,title={Before you move on}]
What does \gu{ગાંઠ} mean? (\textquotedblleft{}A lump; a knot\textquotedblright{}.) Say it once more.
\end{tcolorbox}
\section[\texorpdfstring{\gu{પથારીવશ} (pathārīvaś)}{પથારીવશ (pathārīvaś)}]{\gu{પથારીવશ} (pathārīvaś) — bedridden}

\label{lesson:GU-C479-patharivas}



\begin{quote}
Before the new one: say the Gujarati for a convulsion, then the Gujarati for a lump.
\end{quote}

\subsection*{You'll want to know: \gu{પથારીવશ}}
\textbf{\gu{પથારીવશ}} — \emph{pathārīvaś} — \textquotedblleft{}bedridden\textquotedblright{}.

\begin{quote}
\gu{તાવને} \gu{લીધે} \gu{દાદી} \gu{એક} \gu{અઠવાડિયું} \gu{પથારીવશ} \gu{રહ્યાં}.
\end{quote}

\begin{quote}
\emph{tāvne līdhe dādī ek aṭhvāḍiyũ pathārīvaś rahyā̃.}
\end{quote}

\begin{quote}
Because of the fever my grandmother was bedridden for a week.
\end{quote}

\begin{cousinweb}
\textbf{\gu{પથારીવશ}} (\emph{pathārīvaś}) is built from \textbf{\gu{પથારી}} (\emph{pathārī}), \textquotedblleft{}bedding\textquotedblright{}, and \emph{vaś}, from Sanskrit \emph{vaśa}, \textquotedblleft{}will, power\textquotedblright{} (Monier-Williams): in the power of the bed.
\end{cousinweb}

\begin{grammarlens}[title={What you've built}]
One more word for describing what you lived and felt.
\end{grammarlens}

\subsection*{Your turn}
\begin{itemize}
\raggedright
  \item \emph{Say it:} \emph{pathārīvaś}
  \item \emph{Say it:} \emph{pathārīvaś}, once more
  \item \emph{Say it:} \emph{tāvne līdhe dādī ek aṭhvāḍiyũ pathārīvaś rahyā̃.} — Because of the fever my grandmother was bedridden for a week.
  \item \emph{Recall:} say the Gujarati for a convulsion, then the Gujarati for a lump, then say \emph{pathārīvaś} again
\end{itemize}

\begin{tcolorbox}[breakable,colback=teal!4,colframe=teal!35!black,title={Before you move on}]
What does \gu{પથારીવશ} mean? (\textquotedblleft{}Bedridden\textquotedblright{}.) Say it once more.
\end{tcolorbox}
